\documentclass[11pt,a4paper]{article}

% ---------------------------------------------------------
% PREAMBLE (stable, warning-free, Overleaf-safe)
% ---------------------------------------------------------
\usepackage[utf8]{inputenc}
\usepackage[T1]{fontenc}
\usepackage{lmodern}
\usepackage{amsmath, amssymb, amsfonts}
\usepackage{bm}
\usepackage{geometry}
\usepackage{graphicx}
\usepackage{hyperref}
\usepackage{cite}
\usepackage{authblk}
\usepackage{physics}
\usepackage{mathtools}

\geometry{margin=2.4cm}

\title{\textbf{Companion LV: The CMB-Lensing × Galaxy Cross-Correlation in the Relaxation-Driven Cyclic Cosmology}}
\author{Michael Lehmann \\ \texttt{mi.lehmann@gmx.de}}
\date{April 2026}

\begin{document}
\maketitle

\begin{abstract}
The cross-correlation between CMB lensing and galaxy density provides a powerful 
probe of the relationship between matter and the gravitational potential. In the 
standard cosmological model, this correlation is determined by the matter power 
spectrum, the galaxy bias, and the lensing kernel. In the Relaxation-Driven Cyclic 
Cosmology (RDCC), the scale-dependent evolution of the gravitational potential 
modifies the lensing kernel, the matter distribution, and the growth rate in a 
characteristic way. This Companion develops a continuous description of the 
CMB-lensing × galaxy cross-correlation in RDCC, deriving how the relaxation scale 
influences the coupling between matter and lensing, the coherence of large-scale 
modes, and the observational signatures in galaxy surveys and CMB experiments. The 
resulting framework provides a distinctive observational signature of the RDCC 
gravitational field.
\end{abstract}
% ---------------------------------------------------------
\section{Introduction}

The cross-correlation between CMB lensing and galaxy density probes the connection 
between the gravitational potential and the matter distribution. CMB lensing traces 
the integrated potential along the line of sight, while galaxy surveys trace the 
matter density at specific redshifts. Their cross-correlation therefore encodes the 
relationship between matter and the gravitational field.

In the standard cosmological model, this correlation is determined by the matter 
power spectrum, the galaxy bias, and the lensing kernel. In RDCC, the gravitational 
potential evolves differently. The relaxation mechanism suppresses the decay of 
small-scale modes and enhances the coherence of large-scale modes. This modifies 
the lensing kernel, the matter distribution, and the growth rate in a scale-
dependent manner, producing a distinctive signature in the CMB-lensing × galaxy 
cross-correlation.
% ---------------------------------------------------------
\section{The RDCC Lensing Kernel}

The CMB lensing potential is given by
\begin{equation}
    \phi(\hat{\mathbf{n}})
    = -2 \int_{0}^{\chi_{*}}
      d\chi\,
      \frac{\chi_{*} - \chi}{\chi_{*}\chi}\,
      \left[\Phi + \Psi\right](\chi,\hat{\mathbf{n}}).
\end{equation}

In RDCC, the metric potentials evolve as
\begin{align}
    \Phi_{\mathrm{RDCC}}(k,a)
    &= \Phi(k,a)
       \left[
         1 - \chi_{\Phi}
         \left(\frac{k}{k_{\mathrm{rel}}}\right)^{\alpha}
       \right], \\
    \Psi_{\mathrm{RDCC}}(k,a)
    &= \Psi(k,a)
       \left[
         1 - \chi_{\Psi}
         \left(\frac{k}{k_{\mathrm{rel}}}\right)^{\alpha}
       \right].
\end{align}

The RDCC lensing kernel becomes
\begin{equation}
    W_{\phi,\mathrm{RDCC}}(\chi)
    = W_{\phi}(\chi)
      \left[
        1 - \chi_{\phi}
        \left(\frac{k}{k_{\mathrm{rel}}}\right)^{\alpha}
      \right].
\end{equation}
% ---------------------------------------------------------
\section{The RDCC Galaxy Density Field}

The galaxy density contrast is given by
\begin{equation}
    \delta_{g}(k,z)
    = b(z)\, \delta_{m}(k,z),
\end{equation}
where $b(z)$ is the galaxy bias.

In RDCC, the matter density evolves as
\begin{equation}
    \delta_{m,\mathrm{RDCC}}(k,z)
    = \delta_{m}(k,z)
      \left[
        1 - \chi_{\delta}
        \left(\frac{k}{k_{\mathrm{rel}}}\right)^{\alpha}
      \right].
\end{equation}

The galaxy density becomes
\begin{equation}
    \delta_{g,\mathrm{RDCC}}(k,z)
    = b(z)\,
      \delta_{m}(k,z)
      \left[
        1 - \chi_{g}
        \left(\frac{k}{k_{\mathrm{rel}}}\right)^{\alpha}
      \right].
\end{equation}
% ---------------------------------------------------------
\section{The RDCC Lensing × Galaxy Cross-Correlation}

The cross-correlation is defined as
\begin{equation}
    C_{L}^{\phi g}
    = \int d\chi\,
      \frac{W_{\phi}(\chi)\, W_{g}(\chi)}{\chi^{2}}\,
      P_{m}\left(k=\frac{L}{\chi},z(\chi)\right).
\end{equation}

In RDCC, this becomes
\begin{equation}
    C_{L}^{\phi g,\mathrm{RDCC}}
    = C_{L}^{\phi g}
      \left[
        1 - \chi_{\phi g}
        \left(\frac{L}{L_{\mathrm{rel}}}\right)^{\alpha}
      \right].
\end{equation}

This modification reflects the scale-dependent evolution of the gravitational 
potential and the matter distribution.
% ---------------------------------------------------------
\section{Large-Scale Coherence and RDCC Signatures}

The relaxation mechanism enhances the coherence of large-scale modes. The 
CMB-lensing × galaxy cross-correlation therefore exhibits:

\begin{itemize}
    \item enhanced power at large angular scales,
    \item suppressed power at small scales,
    \item a characteristic turnover near $L_{\mathrm{rel}}$,
    \item a modified redshift dependence of the correlation amplitude.
\end{itemize}

These signatures provide a direct observational probe of the RDCC gravitational 
field.
% ---------------------------------------------------------
\section{Conclusion}

The CMB-lensing × galaxy cross-correlation in the Relaxation-Driven Cyclic 
Cosmology reflects the scale-dependent evolution of the gravitational potential and 
the matter distribution. The relaxation mechanism modifies the lensing kernel, the 
matter density field, and the growth rate in a scale-dependent manner, producing 
distinctive signatures in the cross-correlation spectrum. These effects provide a 
direct observational window into the relaxation scale and the temporal structure of 
the RDCC gravitational field. The present Companion establishes the theoretical 
foundation for future studies of CMB–LSS coupling, matter evolution, and the 
interplay between light and gravity in RDCC.

\bibliographystyle{apsrev4-2}
\nocite{*}
\bibliography{references}

\end{document}
